Predicting three-dimensional grain growth requires connecting local grain-boundary evolution to long-time changes in grain size, topology, and morphology. Physics-based simulations can generate the large datasets needed to train machine-learning surrogates, but a surrogate inherits both the intended dynamics and the numerical artifacts of its training source. This dissertation develops a source-aware local-learning framework that progresses from physically controlled stochastic simulations to direct learning from sparse four-dimensional experiments.

The first completed foundation establishes that the geometry and stochastic sampling of a local simulation neighborhood control lattice pinning and inclination-dependent anisotropy. Gaussian sampling suppresses strong lattice-aligned preferences, deliberately reshaped neighborhoods prescribe graded anisotropy, and rotation of the neighborhood rotates the resulting inclination response. These results show that the simulation neighborhood is part of the physical content encoded in training labels, although they do not establish universal material fidelity~\cite{tian2026neighborhood}.

The second completed foundation, 3D-PRIMME, demonstrates that a local neural operator can learn controlled mode-filter evolution from minimal temporal supervision. The model reproduces major kinetic and topological trends, learns prescribed inclination dependence, and transfers without retraining from \(100^3\)-voxel training volumes to domains as large as \(1024^3\). However, this performance represents fidelity to the simulation teacher rather than proof of universal grain-growth physics~\cite{tian2026threeDprimme}.

Ongoing work extends the same architecture to a five-state laboratory diffraction-contrast-tomography dataset. Stable preliminary results demonstrate residual-driven window selection, integer registration without label interpolation, grain linkage, direct training from one experimental interval, replicate rollouts, held-out evaluation, and substantial suppression of artificial drift. Early coarsening and normalized grain-size distributions are reproduced, while late-time slowdown, topology, and some individual-grain responses remain unresolved~\cite{tian2026experiment}.

The proposed research has three aims. Aim 1 will establish reliability across multiple experimental windows and quantify empirical variability due to data partition, model initialization, and key curation choices. Aim 2 will perform two paired validations: simulation reference versus simulation-trained prediction, and experimental reference versus experiment-trained prediction, using the metrics already established in the corresponding manuscripts and adding further metrics as needed. Aim 3 will combine input-side scale attribution with output-side response and Jacobian analysis to determine which spatial information a reliable experiment-trained model uses and how that information is converted into local evolution decisions. Together, the research will establish when a local learned operator can be trusted relative to its data source and create a defensible basis for physical interpretation of measured mesoscale evolution.
